\section{Экспериментальная апробация прототипа на сценариях регрессий}
\label{sec:razdel5}

\subsection{Сценарии внесения регрессий и диагностика через pprof-профиль}
\label{sec:5.1}

Для апробации прототипа подготовлены четыре независимые ветки
\texttt{regression/\allowbreak *}, каждая из которых представляет собой один точечный
diff от \texttt{main}, вносящий один типичный класс регрессии
производительности (таблица~\ref{tab:5.1}), а также контрольный сценарий —
\texttt{main} без изменений, используемый для оценки доли ложных
срабатываний (раздел~\ref{sec:5.2}).

\begin{table}[htbp]
\caption{Сценарии регрессий, апробированные на прототипе}
\label{tab:5.1}
\centering
\begin{tabular}{p{3.4cm}p{5.5cm}p{4.8cm}}
\toprule
Ветка & Изменение & Затрагиваемая метрика \\
\midrule
\texttt{regression/\allowbreak extra-\allowbreak allocation} & убран \texttt{capacity hint} у \texttt{make([]Record, 0, len(records))} в \texttt{FilterActive} & B/op, allocs/op (рост) \\
\texttt{regression/\allowbreak quadratic-\allowbreak dedup} & хеш-таблица в \texttt{Deduplicate} заменена вложенным линейным поиском & ns/op (рост); см. ниже — B/op при этом \emph{падает} \\
\texttt{regression/\allowbreak latency} & добавлен \texttt{time.Sleep(1μs)} в горячий путь \texttt{FindByID} & ns/op (рост) \\
\texttt{regression/\allowbreak memory-\allowbreak growth} & добавлен никогда не очищаемый пакетный кэш \texttt{aggregateHistory} в \texttt{Aggregate} & B/op, allocs/op (рост, неограниченно) \\
\bottomrule
\end{tabular}
\end{table}

Каждый сценарий выбран так, чтобы задействовать разный механизм деградации:
явный рост аллокаций, изменение класса алгоритмической сложности
($O(n) \to O(n^2)$, формализовано в разделе~\ref{sec:2.2}), синхронная
задержка на горячем пути и неограниченный рост занимаемой памяти — этот
последний тип регрессии на фиксированном единичном прогоне бенчмарка не
отличим от простого роста аллокаций, но исторически является одной из
наиболее опасных категорий продакшен-инцидентов (утечка памяти,
проявляющаяся только при длительной работе процесса).

\textbf{Диагностическая ценность pprof-профиля.} Наиболее показательный
результат апробации получен на сценарии \texttt{regression/\allowbreak quadratic-dedup}
и напрямую подтверждает вывод раздела~\ref{sec:1.2} о необходимости
диагностического pprof-шага в дополнение к агрегированному сравнению.
Прогон \texttt{scripts/\allowbreak profile-diff.sh} (\texttt{COUNT=10}) на этой ветке
против эталона дал:

\begin{itemize}
\item \textbf{CPU-профиль} (\texttt{-diff\_base}, cumulative time):
функция \texttt{Deduplicate} получает $\delta_f = {+6{,}72}$с
\texttt{flat}-времени (около 7,8\% от общего числа сэмплов профиля) — рост
ожидаем, поскольку вложенный цикл поиска теперь исполняется непосредственно
в теле функции, а не делегируется внутренним примитивам map в рантайме.
Одновременно функции \texttt{runtime.mapaccess2\_fast64} и
\texttt{runtime.mapassign\_fast64} получают отрицательную $\delta_f$
(они больше не вызываются вовсе, поскольку map \texttt{seen} убрана из
кода).
\item \textbf{Heap-профиль} (\texttt{-diff\_base}, \texttt{alloc\_space}):
та же функция \texttt{Deduplicate} получает $\delta_f \approx
{-19\,674}$МБ — то есть регрессия \texttt{quadratic-dedup}
\emph{уменьшает} суммарный объём выделенной памяти, поскольку вместе с
map \texttt{seen} исчезла и её аллокация.
\end{itemize}

Это контринтуитивный, но воспроизводимый результат: регрессия является
чистой CPU-регрессией (рост ns/op при неизменной или даже улучшенной
метрике allocs/op и B/op). Если бы критерий~\eqref{eq:criterion}
применялся только к метрикам аллокации (что является частой практикой при
беглом код-ревью — "визуально страшнее" изменение map на цикл не всегда
ассоциируется с ростом именно времени), эта регрессия могла бы быть
пропущена или даже неверно истолкована как улучшение по B/op. Обнаружить
и объяснить её позволяет именно сочетание метрики ns/op (даёт
\texttt{benchstat}) и функционального CPU-профиля (даёт
\texttt{profile-diff.sh}) — практическое подтверждение архитектурного
решения раздела~\ref{sec:3.2} использовать оба механизма параллельно, а не
взаимоисключающе.

\subsection{Выявление отклонений и проверка корректности работы системы}
\label{sec:5.2}

Для проверки корректности критерия~\eqref{eq:criterion} выполнен
\texttt{scripts/\allowbreak experiments/\allowbreak sensitivity.sh}: для каждого из пяти сценариев
(контрольный \texttt{main} и четыре ветки \texttt{regression/\allowbreak *}) снят один
прогон бенчмарков (\texttt{-count=10}), после чего результат прогнан через
\texttt{benchstat} и критерий~\eqref{eq:criterion} по всем сочетаниям сетки
$\alpha \in \{0{,}01;\ 0{,}05;\ 0{,}10\}$, $T \in \{5;10;15;20;30;50\}\%$ —
всего 90 комбинаций (\texttt{experiments/\allowbreak sensitivity\_results.csv}).
Прогон выполнен на раннере \texttt{ubuntu-latest} GitHub Actions
(\texttt{workflow\_dispatch} \texttt{experiments.yml}), а не локально: эталон
(\texttt{benchmarks/\allowbreak baseline.txt}) специально пересобран на том же классе
раннера непосредственно перед экспериментом, чтобы сравнение "кандидат vs.
эталон" не смешивалось с разницей в железе между эталоном и кандидатом
(раздел~\ref{sec:razdel4}).

\begin{table}[htbp]
\caption{Обнаружение регрессии в зависимости от порога $T$ и уровня значимости $\alpha$ (раннер GitHub Actions, \texttt{experiments/sensitivity\_results.csv})}
\label{tab:5.2}
\centering
\begin{tabular}{llcccccc}
\toprule
Сценарий & $\alpha$ & $T{=}5\%$ & $T{=}10\%$ & $T{=}15\%$ & $T{=}20\%$ & $T{=}30\%$ & $T{=}50\%$ \\
\midrule
\multirow{3}{*}{\texttt{main} (без регрессии)}
 & 0,01 & 0 & 0 & 0 & 0 & 0 & 0 \\
 & 0,05 & 0 & 0 & 0 & 0 & 0 & 0 \\
 & 0,10 & \textbf{1} & \textbf{1} & \textbf{1} & 0 & 0 & 0 \\
\midrule
\texttt{extra-allocation} & любое & 1 & 1 & 1 & 1 & 1 & 1 \\
\texttt{quadratic-dedup} & любое & 1 & 1 & 1 & 1 & 1 & 1 \\
\texttt{latency} & любое & 1 & 1 & 1 & 1 & 1 & 1 \\
\texttt{memory-growth} & любое & 1 & 1 & 1 & 1 & 1 & 1 \\
\bottomrule
\end{tabular}
\end{table}

Таблица~\ref{tab:5.2} (1 = регрессия обнаружена хотя бы по одной метрике)
даёт две устойчивые закономерности, обе — на реальном раннере GitHub Actions
заметно чище, чем на предварительном локальном прогоне на ноутбуке
разработчика. Во-первых, все четыре внесённые регрессии обнаруживаются во
\emph{всём} исследованном диапазоне порогов, вплоть до $T=50\%$, при
\emph{самом строгом} из исследованных уровней значимости $\alpha=0{,}01$ —
относительное отклонение $\Delta_{b,m}$, вносимое каждым из четырёх
сценариев, значительно превышает 50\% хотя бы по одной метрике при высокой
статистической значимости, то есть регрессии "с запасом" отделимы от шума.
Во-вторых, на контрольном сценарии (\texttt{main} против самого себя, без
внесённых изменений) при $\alpha \le 0{,}05$ ложных срабатываний не
возникает \emph{ни при одном} из исследованных порогов, вплоть до
$T=5\%$ — то есть на этом раннере одной лишь статистической значимости
критерия (без учёта практического порога) уже достаточно, чтобы отфильтровать
шум окружения. Ложные срабатывания появляются только при самом мягком
$\alpha=0{,}10$ и только при $T \le 15\%$ — граница ложных срабатываний
лежит между $T=15\%$ и $T=20\%$ для этого уровня значимости. Эмпирическая
рекомендация по результатам эксперимента: $T \ge 20\%$ достаточен даже при
мягком $\alpha=0{,}10$, а при $\alpha \le 0{,}05$ пороговое условие можно не
завышать вовсе — что подтверждает необходимость сочетания условий из
формулы~\eqref{eq:criterion}: на разных уровнях значимости именно их
комбинация, а не любое из условий по отдельности, определяет границу
ложных срабатываний.

Дополнительно, \texttt{scripts/\allowbreak experiments/\allowbreak power\_analysis.sh}
(\texttt{experiments/\allowbreak power\_results.csv}) проверяет устойчивость решения
при варьировании числа повторов \texttt{-count} через НЕЗАВИСИМЫЕ повторные
прогоны (в отличие от \texttt{sensitivity.sh}, где сравнение по сетке
$\alpha \times T$ выполняется на одном и том же прогоне бенчмарков). На том же
раннере выполнено \texttt{TRIALS=5} независимых прогонов для каждого из
$\texttt{count} \in \{3;5;10;20\}$ на сценариях \texttt{main} и
\texttt{quadratic-dedup} (фиксированные $\alpha=0{,}05$, $T=10\%$) — всего 40
независимых запусков \texttt{go test -bench}. На \texttt{quadratic-dedup}
регрессия обнаружена во \emph{всех} 20 из 20 прогонов, включая наименьшее
$\texttt{count}=3$ — мощность критерия для регрессии такой величины
оказалась полной уже при минимальном числе повторов. На контрольном
\texttt{main} из 20 независимых прогонов ложное срабатывание произошло в 2
(при $\texttt{count}=5$, испытание 1, и при $\texttt{count}=10$, испытание
5) — эмпирическая частота ложных срабатываний $\approx 10\%$ при этой
комбинации $\alpha=0{,}05$, $T=10\%$. Показательно, что это \emph{отличается}
от результата \texttt{sensitivity.sh} на том же сочетании $\alpha=0{,}05$,
$T=10\%$, где на \emph{единственном} прогоне \texttt{main} ложных
срабатываний не было (табл.~\ref{tab:5.2}) — расхождение единичного
измерения и результата по 20 независимым прогонам является прямой
эмпирической иллюстрацией тезиса Georges~et~al.~\cite{georges2007} и
Mytkowicz~et~al.~\cite{mytkowicz2009} (раздел~\ref{sec:1.1}) о том, что
единичное измерение производительности систематически вводит в
заблуждение и не заменяет оценку по выборке. Выполненный масштаб
(\texttt{TRIALS=5}, 4 значения \texttt{count}, 2 сценария) существенно
превосходит предварительный локальный прогон (2 испытания на точку), но
всё ещё меньше предусмотренного скриптом предельного масштаба
(\texttt{TRIALS=20}, \texttt{COUNTS} до 50) — дальнейшее увеличение
выборки обозначено как направление продолжения эксперимента в
разделе~\ref{sec:5.4}.

\subsection{Оценка накладных расходов профилирования в CI/CD}
\label{sec:5.3}

Ключевой вопрос практического применения прототипа — во сколько раз (или
на сколько) сбор pprof-профилей замедляет CI-джобу по сравнению с
"обычным" запуском бенчмарков без профилирования, и какой объём
дополнительного места на диске/в артефактах CI это добавляет. Для ответа
на том же раннере GitHub Actions выполнен
\texttt{scripts/\allowbreak experiments/\allowbreak overhead.sh}: 10 независимых прогонов
полного набора бенчмарков (\texttt{-count=10}) в режиме \texttt{plain} (без
флагов профилирования) и 10 — в режиме \texttt{profiled}
(\texttt{-cpuprofile}/\texttt{-memprofile}), с измерением настенного
времени выполнения и суммарного размера полученных profile-файлов
(таблица~\ref{tab:5.3}).

\begin{table}[htbp]
\caption{Настенное время выполнения бенчмарков с профилированием и без, раннер GitHub Actions (experiments/overhead\_results.csv)}
\label{tab:5.3}
\centering
\begin{tabular}{cccc}
\toprule
Прогон & \texttt{plain}, с & \texttt{profiled}, с & Размер профилей, байт \\
\midrule
1  & 80,373 & 83,964 & 66\,823 \\
2  & 80,537 & 83,863 & 65\,445 \\
3  & 80,422 & 83,654 & 66\,593 \\
4  & 81,036 & 84,065 & 65\,598 \\
5  & 80,638 & 85,157 & 67\,816 \\
6  & 79,406 & 83,554 & 65\,519 \\
7  & 79,869 & 83,860 & 66\,076 \\
8  & 80,219 & 84,068 & 66\,210 \\
9  & 80,804 & 83,462 & 65\,291 \\
10 & 80,224 & 84,057 & 65\,388 \\
\midrule
Среднее & 80,35 & 83,97 & 66\,076 ($\approx$66,1 КБ) \\
Ст.\ откл. & 0,47 & 0,47 & --- \\
\bottomrule
\end{tabular}
\end{table}

\begin{figure}[htbp]
\centering
\begin{tikzpicture}
\begin{axis}[
  ybar,
  bar width=28pt,
  symbolic x coords={plain,profiled},
  xtick=data,
  ylabel={Настенное время, с},
  width=0.62\textwidth,
  height=6cm,
  ymin=0,
  nodes near coords,
  nodes near coords align={vertical},
  error bars/y dir=both,
  error bars/y explicit,
]
\addplot+[black, fill=gray!35] coordinates {
  (plain,80.35) +- (0,0.47)
  (profiled,83.97) +- (0,0.47)
};
\end{axis}
\end{tikzpicture}
\caption{Среднее настенное время выполнения набора бенчмарков: без профилирования и с профилированием (столбики погрешности --- ст.\ откл.\ по 10 прогонам, раннер GitHub Actions)}
\label{fig:5.1}
\end{figure}

Относительное увеличение среднего настенного времени составило
$(83{,}97 - 80{,}35)/80{,}35 \times 100\% \approx 4{,}5\%$. В отличие от
предварительного локального прогона (5 испытаний на режим на ноутбуке
разработчика, где разница средних не превышала межпрогонный шум), на
раннере GitHub Actions межпрогонный разброс оказался заметно меньше
(стандартное отклонение $\approx$0,47с в обоих режимах) при кратно большей
разнице средних (3,62с) — то есть на выборке из 10 прогонов на режим
отличие статистически уверенно отделимо от шума измерения: включение
CPU/heap-профилирования действительно измеримо замедляет прогон
бенчмарков, но на величину $\approx$4,5\%, что для цели данной работы
(обнаружение регрессий с порогом $T$ в единицы--десятки процентов, см.
раздел~\ref{sec:5.2}) остаётся практически приемлемым: постоянное
включение профилирования не искажает решение "регрессия/не регрессия" по
основному (агрегированному) критерию, поскольку сравнение
$\Delta_{b,m}$ в формуле~\eqref{eq:criterion} выполняется между
одинаково профилированными эталоном и кандидатом, а не между
профилированным и непрофилированным прогоном. Второй компонент накладных
расходов — объём артефактов: средний суммарный размер CPU- и
heap-профилей одного прогона составил $\approx$66,1~КБ, что
пренебрежимо мало относительно типичных квот на артефакты CI (сотни
мегабайт -- гигабайты) даже при накоплении профилей за сотни PR.

Выполненный масштаб (10 прогонов на режим на выделенном раннере, а не на
разделяемом ноутбуке) достаточен для уверенного вывода о порядке величины
накладных расходов ($\approx$4--5\%) и о том, что это отличие не является
шумом измерения — стандартное отклонение на раннере оказалось на порядок
меньше, чем в предварительном локальном прогоне, что само по себе
иллюстрирует важность выполнения подобных измерений на выделенном,
однородном оборудовании (раздел~\ref{sec:razdel4}), а не на разделяемой
рабочей машине.

\subsection{Выводы}
\label{sec:5.4}

\begin{enumerate}
\item Апробированы четыре сценария регрессии производительности,
затрагивающие разные механизмы деградации (рост аллокаций, изменение
класса алгоритмической сложности, синхронная задержка, неограниченный
рост памяти); на реальном раннере GitHub Actions все четыре обнаружены
критерием~\eqref{eq:criterion} во \emph{всём} исследованном диапазоне
порогов $T \in [5\%, 50\%]$ при самом строгом уровне значимости
$\alpha=0{,}01$ — регрессии отделимы от шума с большим запасом.
\item На сценарии \texttt{regression/\allowbreak quadratic-dedup} экспериментально
показано, что изменение алгоритмической сложности может \emph{одновременно}
увеличивать время выполнения (ns/op, от $+62\%$ до $+69\%$ в независимых
прогонах, см. рис.~\ref{fig:4.1}) и уменьшать объём аллокаций (B/op,
$\delta \approx -19\,674$МБ по heap-профилю) — то есть
регрессия, невидимая при слежении только за метриками памяти, надёжно
обнаруживается через ns/op и локализуется до конкретной функции через
pprof CPU-профиль, что подтверждает необходимость двухуровневой схемы
профилирования, спроектированной в разделе~\ref{sec:3.2}.
\item На контрольном сценарии (\texttt{main} без изменений) при
$\alpha \le 0{,}05$ ложных срабатываний не возникло ни при одном из
исследованных порогов вплоть до $T=5\%$; при $\alpha=0{,}10$ граница
ложных срабатываний лежит между $T=15\%$ и $T=20\%$ — практическая
рекомендация по результатам эксперимента: $T \ge 20\%$ достаточен даже
при мягком уровне значимости, а при $\alpha \le 0{,}05$ порог можно не
завышать. Это заметно чище границы, полученной в предварительном
локальном прогоне на разделяемом ноутбуке ($T \ge 30\%$), что подтверждает
значимость пересборки эталона на однородном выделенном раннере
(раздел~\ref{sec:razdel4}).
\item Независимые повторные прогоны (\texttt{power\_analysis.sh}, 40
запусков) показали полную мощность критерия на \texttt{quadratic-dedup}
(20 из 20, включая минимальное $\texttt{count}=3$) и эмпирическую частоту
ложных срабатываний $\approx$10\% на \texttt{main} (2 из 20 при
$\alpha=0{,}05$, $T=10\%$) — при том, что единственный прогон
\texttt{sensitivity.sh} на той же комбинации параметров ложных
срабатываний не показал. Это расхождение единичного измерения и оценки по
20 независимым прогонам — прямая эмпирическая иллюстрация тезиса
Georges~et~al.~\cite{georges2007} и Mytkowicz~et~al.~\cite{mytkowicz2009}
о недостаточности единичного измерения производительности.
\item Оценены накладные расходы профилирования в CI/CD на выделенном
раннере (10 прогонов на режим): среднее увеличение настенного времени
выполнения бенчмарков составило $\approx$4,5\%, что на этот раз
статистически уверенно отличимо от межпрогонного шума измерения
(стандартное отклонение $\approx$0,47с против разности средних 3,62с);
дополнительный объём артефактов на один прогон составил
$\approx$66,1~КБ. Оба показателя подтверждают, что постоянное включение
pprof-профилирования в каждый прогон CI практически не увеличивает
стоимость пайплайна, при этом само отличие от предварительной локальной
оценки ($\approx$2,2\%, статистически неразличимой от шума) показывает
ценность выполнения подобных измерений на выделенном, а не разделяемом
оборудовании.
\end{enumerate}
